% !TeX root = ../navier-stokes-manuscript.tex
% ============================================================
% 2.1 Vortex Stretching
% ============================================================

The momentum equation has four pieces:
\[
  \underbrace{\partial_t u}_{\text{local acceleration}}
  +
  \underbrace{(u\cdot\nabla)u}_{\text{nonlinear transport}}
  =
  \underbrace{-\nabla p}_{\text{pressure}}
  +
  \underbrace{\nu\Delta u}_{\text{viscous diffusion}}.
\]
The central regularity problem lies in the competition between nonlinear transport and viscosity.
The hope is simple: viscosity smooths the fluid faster than nonlinear transport can sharpen it.
If that remains true at every scale, the velocity stays smooth.

\subsection{Vortex Stretching}\label{sub:ThePerfectlyStretchingVortex}

One narrow material vortex tube puts that hope under pressure immediately.
Along a material segment of its axis, the surrounding flow can pull the two ends apart and lengthen the rotating core.
The part of the velocity gradient that records this deformation is the symmetric strain
\[
  S:=\frac12\bigl(\nabla u+(\nabla u)^{\mathsf T}\bigr),
\]
and if \(e\) points along the segment while \(L(t)\) denotes its length, extension along \(e\) at rate \(\sigma(t)>0\) means
\[
  Se=\sigma(t)e,
  \qquad
  \sigma(t)>0,
  \qquad
  \frac{\dot L(t)}{L(t)}
  =
  e\cdot Se
  =
  \sigma(t).
\]
The segment is being stretched along its own axis.

Because the fluid is incompressible, that same motion squeezes the cross-section.
Let \(\mathcal V\) be the fixed material volume, set \(A(t):=\mathcal V/L(t)\), and write \(r_{\mathrm{eff}}(t):=\sqrt{A(t)/\pi}\).
Then
\[
  \nabla\cdot u=0,
  \qquad
  \frac{d}{dt}(A L)=0,
  \qquad
  \frac{\dot A}{A}=-\sigma(t),
  \qquad
  \frac{\dot r_{\mathrm{eff}}}{r_{\mathrm{eff}}}
  =
  -\frac{\sigma(t)}{2}.
\]
Lengthening and narrowing are not two separate events.
The same strain that increases \(L(t)\) drives the effective radius down.

That narrowing acts on the tube's rotation.
With \(\omega:=\nabla\times u\), the vorticity satisfies along the moving fluid
\[
  \frac{D\omega}{Dt}
  =
  S\omega+\nu\Delta\omega,
  \qquad
  \frac{D}{Dt}
  :=
  \partial_t+u\cdot\nabla,
\]
and when the vorticity aligns with the stretching direction, \(\omega=|\omega|e\),
\[
  S\omega
  =
  \sigma(t)\omega,
  \qquad
  \frac12\frac{D|\omega|^2}{Dt}
  =
  \sigma(t)|\omega|^2
  +
  \nu\,\omega\cdot\Delta\omega.
\]
Axial stretching now appears twice in the same material piece: it lengthens the segment and multiplies aligned vorticity.
The viscous term remains in that balance, with no pointwise sign forced by the identity itself.
So amplification only occurs on intervals where the full right-hand side stays positive.

What grows here is the rotating part of the velocity gradient, not the kinetic energy of the segment as a whole.
If the speed stays bounded by \(U_0\) on the fixed-volume material segment \(\Omega_{\mathrm{seg}}(t)\), then
\[
  E_{\mathrm{seg}}(t)
  :=
  \frac12\int_{\Omega_{\mathrm{seg}}(t)}|u(x,t)|^2\,dx
  \leq
  \frac12U_0^2\mathcal V
  <
  \infty,
  \qquad
  \left|\frac12\bigl(\nabla u-(\nabla u)^{\mathsf T}\bigr)\right|_{\mathrm F}
  =
  \frac{|\omega|}{\sqrt2}
  \leq
  |\nabla u|_{\mathrm F}.
\]
A bounded-speed segment of fixed material volume can keep finite kinetic energy while its rotating gradient grows and its core narrows.
The gradients rise.

\divider{\NavierStokes{} Scaling}

That falling width is exactly where the original hope would have to improve.
If viscosity were going to outrun nonlinear sharpening, some relative gain should appear when the same equation is examined at the smaller scale now being created.
The comparison has to be made with a rescaled solution, not by pretending that a later state of this material tube has already been reached.
Fix \(t_0\) and set \(\ell:=r_{\mathrm{eff}}(t_0)\).
For \(\lambda>1\), define
\[
  u_\lambda(x,t)
  :=
  \lambda u(\lambda x,\lambda^2t),
  \qquad
  p_\lambda(x,t)
  :=
  \lambda^2p(\lambda x,\lambda^2t).
\]
At time \(\lambda^{-2}t_0\), the corresponding vortex in \(u_\lambda\) has width \(\lambda^{-1}\ell\).
Differentiating the rescaled fields gives
\[
\begin{aligned}
  \partial_tu_\lambda(x,t)
  &=
  \lambda^3(\partial_tu)(\lambda x,\lambda^2t),
  \\
  (u_\lambda\cdot\nabla)u_\lambda(x,t)
  &=
  \lambda^3\bigl((u\cdot\nabla)u\bigr)(\lambda x,\lambda^2t),
  \\
  \nabla p_\lambda(x,t)
  &=
  \lambda^3(\nabla p)(\lambda x,\lambda^2t),
  \\
  \nu\Delta u_\lambda(x,t)
  &=
  \lambda^3\nu(\Delta u)(\lambda x,\lambda^2t).
\end{aligned}
\]
The narrowed width changes none of the equation's internal proportions.
Local acceleration, nonlinear transport, pressure, and viscosity all scale the same way.
The hoped-for advantage does not appear here.

\divider{Critical Height}

Once the equation stays level under that magnification, the next issue is no longer the equation itself but the ruler being used to measure the field.
A norm can grow or shrink because the scale changed even when the competition did not.
Let \(\Lambda:=(-\Delta)^{1/2}\).
At Sobolev height \(s\), a change of variables gives
\[
  \norm{u_\lambda(t)}_{\dot H^s}^2
  =
  \lambda^{2s-1}
  \norm{u(\lambda^2t)}_{\dot H^s}^2.
\]
The factor \(\lambda^{2s-1}\) is contributed by scaling alone.
It disappears only at \(s=\tfrac12\).
Set
\[
  H(t)
  :=
  \frac12\norm{u(t)}_{\dot H^{1/2}}^2
  =
  \frac12\norm{\Lambda^{1/2}u(t)}_2^2.
\]
At this height the scale change itself is silent, while the levels below and above it still move:
\[
  \norm{u_\lambda(t)}_2^2
  =
  \lambda^{-1}\norm{u(\lambda^2t)}_2^2,
  \qquad
  H_\lambda(t)=H(\lambda^2t),
  \qquad
  \norm{\nabla u_\lambda(t)}_2^2
  =
  \lambda\norm{\nabla u(\lambda^2t)}_2^2.
\]
Kinetic energy drops under concentration, first-derivative energy rises, and the critical height stays fixed.
Finite energy still leaves the concentration problem open.
A change in \(H\) has to come from the equation itself.

\divider{Nonlinear Production versus Viscous Dissipation}

That is where the opening competition takes its exact form.
Apply \(\Lambda^{1/2}\) to the momentum equation and pair with \(\Lambda^{1/2}u\).
Write the signed contribution of nonlinear transport as
\[
  \mathcal P_{1/2}(t)
  :=
  -\operatorname{Re}
  \pair
    {\Lambda^{1/2}u(t)}
    {\Lambda^{1/2}\bigl((u(t)\cdot\nabla)u(t)\bigr)}_{L^2}.
\]
The pressure pairing vanishes because \(\nabla\cdot u=0\), while the Laplacian contributes \(-\nu\norm{\Lambda^{3/2}u}_2^2\).
So
\[
  H'(t)
  +
  \nu\norm{\Lambda^{3/2}u(t)}_2^2
  =
  \mathcal P_{1/2}(t).
\]
The critical height rises only when nonlinear production exceeds viscous dissipation.
Under the same rescaling,
\[
  \mathcal P_{1/2}[u_\lambda](t)
  =
  \lambda^2\mathcal P_{1/2}[u](\lambda^2t),
  \qquad
  \nu\norm{\Lambda^{3/2}u_\lambda(t)}_2^2
  =
  \lambda^2\nu\norm{\Lambda^{3/2}u(\lambda^2t)}_2^2.
\]
The balance remains level here as well.
Nonlinear production and viscous dissipation gain the same \(\lambda^2\) factor.
What changes is the clock attached to the comparison, since the same rescaling reads the field at time \(\lambda^{-2}t\).

\divider{The Heat Clock}

Viscosity already carries that clock inside its linear action.
For the heat equation \(v_t=\nu\Delta v\),
\[
  \widehat v(\xi,t+\delta t)
  =
  e^{-\nu|\xi|^2\delta t}\widehat v(\xi,t).
\]
A structure localized to width \(r\) lives at frequencies \(|\xi|\simeq r^{-1}\).
An order-one viscous change occurs when \(\nu|\xi|^2\delta t\simeq1\), so viscosity associates that width with the time
\[
  \tau_\nu(r)
  :=
  \frac{r^2}{\nu}.
\]
This is not a declared lifetime for the material vortex.
It is the time viscosity alone needs to alter a structure of width \(r\) by order one.
At the finer comparison width,
\[
  \tau_\nu(\lambda^{-1}r)
  =
  \lambda^{-2}\tau_\nu(r),
\]
which is exactly the contraction already built into \(t\mapsto\lambda^{-2}t\).
So every time the width is pushed down by \(\lambda^{-1}\), the viscous comparison time falls by \(\lambda^{-2}\) while the nonlinear--viscous balance itself remains untilted.
The unresolved spatial competition has acquired its temporal image.

\divider{The Zeno Cascade}

Repeat that narrowing geometrically:
\[
  r_n:=2^{-n}r_0,
  \qquad
  \tau_n:=\frac{r_n^2}{\nu}.
\]
Then
\[
  \tau_n
  =
  \frac{r_0^2}{\nu}4^{-n},
  \qquad
  \sum_{n=0}^{\infty}\tau_n
  =
  \frac{4r_0^2}{3\nu}
  <
  \infty.
\]
The clocks shrink fast enough to fit infinitely many finer comparisons into finite total time.
That still does not produce a singularity.
It only describes the schedule that a candidate singular history would have to realize.
One actual solution would have to pass through widths
\[
  r_0>r_1>r_2>\cdots\longrightarrow0
\]
at times
\[
  t_0<t_1<t_2<\cdots\uparrow T_*<\infty,
  \qquad
  t_{n+1}-t_n\asymp\tau_n,
\]
with comparison constants independent of \(n\).
Then the remaining time after stage \(n\) would satisfy
\[
  T_*-t_n
  \asymp
  \sum_{k=n}^{\infty}\tau_k
  =
  \frac{4r_n^2}{3\nu}.
\]
By the time the width has reached \(r_n\), only order \(r_n^2/\nu\) of comparison time remains.
So the Zeno sum leaves finite time available, but only on the condition that one and the same smooth velocity history can keep generating finer concentration on a quadratically collapsing clock.
That is no longer just a spatial schedule.
It has become a demand on temporal response.